% =============================================================================
% Chapter: SQL Basics
% =============================================================================

\chapter{SQL Basics}

\begin{formulabox}[Key SQL Syntax - MEMORIZE!]
\begin{lstlisting}[language=SQL]
SELECT [DISTINCT] column1, column2, aggregate(column)
FROM table1 [JOIN table2 ON condition]
WHERE condition           -- Filter BEFORE grouping
GROUP BY column           -- Group rows
HAVING aggregate_condition -- Filter AFTER grouping
ORDER BY column [ASC|DESC]
LIMIT n;
\end{lstlisting}

\textbf{Execution Order:} FROM → WHERE → GROUP BY → HAVING → SELECT → ORDER BY
\end{formulabox}

% =============================================================================
\section{SELECT Basics}

\subsection{Column Selection}

\begin{examplebox}[Basic SELECT]
\begin{lstlisting}[language=SQL]
-- All columns
SELECT * FROM Employee;

-- Specific columns
SELECT name, salary FROM Employee;

-- With alias
SELECT name AS employee_name, salary * 12 AS annual
FROM Employee;

-- Remove duplicates
SELECT DISTINCT dept FROM Employee;
\end{lstlisting}
\end{examplebox}

\subsection{WHERE Clause}

\begin{center}
\begin{tabular}{ll}
\toprule
\textbf{Operator} & \textbf{Example} \\
\midrule
$=$, $<>$, $<$, $>$ & \texttt{salary > 50000} \\
\texttt{BETWEEN} & \texttt{age BETWEEN 20 AND 30} \\
\texttt{IN} & \texttt{dept IN ('CS', 'Math')} \\
\texttt{LIKE} & \texttt{name LIKE 'A\%'} \\
\texttt{IS NULL} & \texttt{manager IS NULL} \\
\texttt{AND, OR, NOT} & \texttt{dept='CS' AND salary > 50000} \\
\bottomrule
\end{tabular}
\end{center}

\begin{examplebox}[WHERE Examples]
\begin{lstlisting}[language=SQL]
-- Pattern matching
SELECT * FROM Employee
WHERE name LIKE 'A%';  -- Starts with A

-- NULL check (NOT with =!)
SELECT * FROM Employee
WHERE manager IS NULL;

-- Range
SELECT * FROM Employee
WHERE salary BETWEEN 40000 AND 60000;
\end{lstlisting}
\end{examplebox}

% =============================================================================
\section{Aggregate Functions}

\begin{formulabox}[Aggregate Functions]
\begin{center}
\begin{tabular}{lll}
\toprule
\textbf{Function} & \textbf{Description} & \textbf{NULL Handling} \\
\midrule
\texttt{COUNT(*)} & Αριθμός γραμμών & Counts ALL \\
\texttt{COUNT(col)} & Αριθμός non-NULL τιμών & Ignores NULL \\
\texttt{SUM(col)} & Άθροισμα & Ignores NULL \\
\texttt{AVG(col)} & Μέσος όρος & Ignores NULL \\
\texttt{MAX(col)} & Μέγιστο & Ignores NULL \\
\texttt{MIN(col)} & Ελάχιστο & Ignores NULL \\
\bottomrule
\end{tabular}
\end{center}
\end{formulabox}

\begin{warningbox}[COUNT(*) vs COUNT(column)]
\begin{lstlisting}[language=SQL]
-- Table: Employee(name, bonus)
-- Data: ('Alice', 100), ('Bob', NULL), ('Carol', 200)

SELECT COUNT(*) FROM Employee;       -- Returns 3
SELECT COUNT(bonus) FROM Employee;   -- Returns 2 (NULL ignored!)
\end{lstlisting}
\end{warningbox}

% =============================================================================
\section{GROUP BY \& HAVING}

\begin{definitionbox}[GROUP BY vs HAVING]
\begin{itemize}
    \item \textbf{WHERE}: Φιλτράρει γραμμές ΠΡΙΝ το grouping
    \item \textbf{HAVING}: Φιλτράρει groups ΜΕΤΑ το grouping
\end{itemize}
\end{definitionbox}

\begin{examplebox}[GROUP BY Example]
\textbf{Query:} Average salary per department, only for departments with $>$ 2 employees.

\begin{lstlisting}[language=SQL]
SELECT dept, AVG(salary) AS avg_salary, COUNT(*) AS emp_count
FROM Employee
WHERE salary > 30000        -- Filter rows first
GROUP BY dept               -- Then group
HAVING COUNT(*) > 2         -- Filter groups
ORDER BY avg_salary DESC;
\end{lstlisting}
\end{examplebox}

\begin{warningbox}[SELECT with GROUP BY Rule]
Στο SELECT μπορούν ΜΟΝΟ:
\begin{itemize}
    \item Στήλες που είναι στο GROUP BY
    \item Aggregate functions
\end{itemize}

\begin{lstlisting}[language=SQL]
-- WRONG!
SELECT name, dept, AVG(salary) FROM Employee GROUP BY dept;
-- 'name' is not in GROUP BY!

-- CORRECT
SELECT dept, AVG(salary) FROM Employee GROUP BY dept;
\end{lstlisting}
\end{warningbox}

% =============================================================================
\section{Solved Exam Problems}

\subsection{2024-25 Exam Question 13 Style}

\begin{examplebox}[Multi-table Query]
\textbf{Tables:}
\begin{itemize}
    \item \texttt{Student(sid, name, dept)}
    \item \texttt{Course(cid, title, credits)}
    \item \texttt{Enrolls(sid, cid, grade)}
\end{itemize}

\textbf{Query:} Find students with average grade $>$ 8.0 in CS department.

\begin{lstlisting}[language=SQL]
SELECT S.name, AVG(E.grade) AS avg_grade
FROM Student S
JOIN Enrolls E ON S.sid = E.sid
WHERE S.dept = 'CS'
GROUP BY S.sid, S.name
HAVING AVG(E.grade) > 8.0
ORDER BY avg_grade DESC;
\end{lstlisting}

\textbf{Step-by-step:}
\begin{enumerate}
    \item JOIN: Συνδέουμε Student με Enrolls
    \item WHERE: Φιλτράρουμε CS students
    \item GROUP BY: Ομαδοποιούμε ανά student
    \item HAVING: Κρατάμε groups με AVG $>$ 8.0
    \item SELECT: Εμφανίζουμε name και average
\end{enumerate}
\end{examplebox}

\subsection{2024-25 Exam Question 14 Style}

\begin{examplebox}[Subquery Example]
\textbf{Query:} Students who have NOT enrolled in any course.

\begin{lstlisting}[language=SQL]
-- Method 1: NOT IN
SELECT name FROM Student
WHERE sid NOT IN (SELECT DISTINCT sid FROM Enrolls);

-- Method 2: NOT EXISTS
SELECT name FROM Student S
WHERE NOT EXISTS (
    SELECT 1 FROM Enrolls E WHERE E.sid = S.sid
);

-- Method 3: LEFT JOIN + NULL check
SELECT S.name
FROM Student S
LEFT JOIN Enrolls E ON S.sid = E.sid
WHERE E.sid IS NULL;
\end{lstlisting}
\end{examplebox}

\subsection{Finding Maximum with Details}

\begin{examplebox}[Max Salary with Employee Details]
\textbf{Query:} Find employee(s) with highest salary.

\begin{lstlisting}[language=SQL]
-- WRONG: Returns one row even if ties exist
SELECT * FROM Employee ORDER BY salary DESC LIMIT 1;

-- CORRECT: Handles ties
SELECT * FROM Employee
WHERE salary = (SELECT MAX(salary) FROM Employee);
\end{lstlisting}
\end{examplebox}

% =============================================================================
\section{Common Patterns}

\begin{center}
\begin{tikzpicture}
    \matrix (m) [
        matrix of nodes,
        nodes={
            draw=primary,
            minimum width=5cm,
            minimum height=1.5cm,
            anchor=center,
            font=\small,
            align=left
        },
        column 1/.style={nodes={fill=ntua-blue!20, font=\bfseries}},
        row 1/.style={nodes={fill=secondary!30, font=\bfseries}},
        column sep=-\pgflinewidth,
        row sep=-\pgflinewidth
    ] {
        Pattern & SQL Template \\
        Count per group & |[align=left]| \texttt{SELECT g, COUNT(*)}\\
                          \texttt{FROM T GROUP BY g} \\
        Top N per group & |[align=left]| \texttt{WITH ranked AS (...)}\\
                          \texttt{ROW\_NUMBER() OVER} \\
        Self-join & |[align=left]| \texttt{FROM T t1, T t2}\\
                    \texttt{WHERE t1.x = t2.y} \\
        NOT IN pattern & |[align=left]| \texttt{WHERE x NOT IN}\\
                         \texttt{(SELECT x FROM ...)} \\
    };
\end{tikzpicture}
\end{center}

% =============================================================================
\section{Common Mistakes}

\begin{warningbox}[Συχνά Λάθη στις Εξετάσεις]
\begin{enumerate}
    \item \textbf{NULL με = instead of IS}:
    \begin{itemize}
        \item[$\times$] \texttt{WHERE bonus = NULL}
        \item[\checkmark] \texttt{WHERE bonus IS NULL}
    \end{itemize}

    \item \textbf{HAVING χωρίς GROUP BY}:
    \begin{itemize}
        \item HAVING πρέπει να χρησιμοποιείται με aggregate
    \end{itemize}

    \item \textbf{Non-grouped columns στο SELECT}:
    \begin{itemize}
        \item[$\times$] \texttt{SELECT name, SUM(x) GROUP BY dept}
        \item[\checkmark] \texttt{SELECT dept, SUM(x) GROUP BY dept}
    \end{itemize}

    \item \textbf{String comparison}:
    \begin{itemize}
        \item[$\times$] \texttt{WHERE name = "Alice"} (double quotes)
        \item[\checkmark] \texttt{WHERE name = 'Alice'} (single quotes)
    \end{itemize}

    \item \textbf{COUNT με NULL}:
    \begin{itemize}
        \item \texttt{COUNT(*)} counts all rows
        \item \texttt{COUNT(col)} ignores NULL values
    \end{itemize}
\end{enumerate}
\end{warningbox}

% =============================================================================
\section{Quick Reference}

\begin{center}
\begin{tikzpicture}
    \node[draw=ntua-blue, fill=ntua-blue!10, rounded corners=5pt, minimum width=14cm, minimum height=7cm] (main) {};

    \node[anchor=north west, font=\bfseries\Large\color{ntua-blue}] at ([xshift=5pt, yshift=-5pt]main.north west) {SQL Execution Order};

    \node[anchor=north west, align=left] at ([xshift=10pt, yshift=-35pt]main.north west) {
        \begin{tikzpicture}
            \node[draw=secondary, fill=secondary!20, rounded corners=3pt, minimum width=2cm] (from) at (0,0) {1. FROM};
            \node[draw=secondary, fill=secondary!20, rounded corners=3pt, minimum width=2cm] (where) at (2.5,0) {2. WHERE};
            \node[draw=secondary, fill=secondary!20, rounded corners=3pt, minimum width=2.5cm] (group) at (5.5,0) {3. GROUP BY};
            \node[draw=secondary, fill=secondary!20, rounded corners=3pt, minimum width=2cm] (having) at (8.5,0) {4. HAVING};
            \node[draw=secondary, fill=secondary!20, rounded corners=3pt, minimum width=2cm] (select) at (0,-1.2) {5. SELECT};
            \node[draw=secondary, fill=secondary!20, rounded corners=3pt, minimum width=2.5cm] (order) at (3,-1.2) {6. ORDER BY};

            \draw[->, thick] (from) -- (where);
            \draw[->, thick] (where) -- (group);
            \draw[->, thick] (group) -- (having);
            \draw[->, thick] (having) -- (10,0) -- (10,-1.2) -- (select.east);
            \draw[->, thick] (select) -- (order);
        \end{tikzpicture}
    };

    \node[draw=success, fill=box-green, rounded corners=5pt, minimum width=6cm, minimum height=2cm]
        at ([xshift=-3cm, yshift=1.5cm]main.south) {};
    \node[align=center, font=\small] at ([xshift=-3cm, yshift=1.5cm]main.south) {
        \textbf{Key Pattern:}\\
        \texttt{SELECT cols, AGG()}\\
        \texttt{FROM T GROUP BY cols}\\
        \texttt{HAVING AGG() > val}
    };

    \node[draw=accent, fill=box-red, rounded corners=5pt, minimum width=5cm, minimum height=2cm]
        at ([xshift=3.5cm, yshift=1.5cm]main.south) {};
    \node[align=center, font=\small] at ([xshift=3.5cm, yshift=1.5cm]main.south) {
        \textbf{Remember:}\\
        WHERE: before grouping\\
        HAVING: after grouping
    };
\end{tikzpicture}
\end{center}
